\documentclass[10pt]{article}
\usepackage{template/template}
\usepackage{amsmath}
\NewsSetKeys{date={2026年11月9日},vol={027},title={每日有趣新闻},author={编辑部},note={待补编短刊：仅刊载已核查材料，尚未完成正式发行。}}
\begin{document}
\begin{documents}
\newsitem{霍尔木兹海峡：谈判信号与航运风险同时存在}{国际政治｜11月8日｜路透社}{美国总统特朗普表示，倘若美伊谈判继续取得进展，再次对伊朗实施打击的可能性很低。不过，这属于有条件的政治表态，并非停火协议，更不能视为海峡通航恢复。路透社同时报道，海峡附近仍有船只遇袭及美伊双方关于海上行动的不同说法；部分爆炸报道未获独立确认。国际油价在新一周开盘时下跌，但价格反应只反映市场对未来风险的重新定价，不证明局势已经缓和。美国与伊朗对于制裁、核问题和赔偿等议题仍存在实质分歧。报道应继续追踪谈判文本、实际航运数据及独立核实的事件记录，而不是仅凭一方的声明推断结果。资料：路透社2026年11月8日报道\cite{reutersIran,reutersOil}。}
\newsitem{澳大利亚央行：AI投资已产生需求，生产率红利尚待证明}{经济与AI｜11月8日｜路透社}{澳大利亚储备银行助理行长Sarah Hunter表示，目前缺乏足够证据证明人工智能已显著提升整体生产率，或导致大规模就业损失。数据中心建设先增加电力、设备与施工需求，在供给约束下可能推高价格；只有当AI被广泛部署、工作流程真正改变后，生产率改善才可能抵消这些压力。她同时指出，部分更易受AI影响的岗位已有招聘和青年就业走弱迹象，但这不足以直接推出整体失业潮。宏观经济上，应区分资本开支形成的当期需求与未来技术扩散带来的潜在供给提升。资料：路透社2026年11月8日报道\cite{reutersRBA}。}
\newsitem{以色列组阁：取得总统授权不等于拥有议会多数}{国际政治｜11月8日｜路透社}{路透社11月8日新闻综述指出，以色列总统选定加迪·艾森科特尝试组织政府。组阁权是一项程序性机会，最终仍取决于政党谈判、议会支持和联盟协议。不同党派对合作对象的否决条件可能令名义上的反对派席位难以自动转化为稳定执政联盟。因此，在实际协议达成前，不宜将组阁尝试写作政权更替已经完成。资料：路透社11月8日新闻综述\cite{reutersIsrael}。}
\newsitem{数学检索记录}{数学｜11月9日}{已核查当日数学新预印本检索入口，但本次尚未获得可独立核实的原始论文作者、定理全文与版本信息，故不刊登未经确认的突破性主张。后续需覆盖数论、代数、几何、组合、逻辑、概率及Lean/mathlib，并逐项核对arXiv原文和同行评议状态。}
\newsitem{编辑说明}{发行状态}{本稿是可追溯来源的短刊草稿，尚未达到深稿篇幅、新闻图片选配、XeLaTeX构建及PDF视觉验收要求；不可作为正式发行版。}
\bibliographystyle{plain}
\bibliography{refs}
\end{documents}
\end{document}
